% ============================================================
\chapter{SQL --- Basic Queries}
\label{ch:sql_basics}
% ============================================================

In 1974, Donald Chamberlin and Raymond Boyce, researchers at IBM, created SEQUEL (\emph{Structured English Query Language}), later renamed SQL for legal reasons. Their ambition: to allow non-programmers to query databases by writing sentences close to English. The idea was revolutionary: instead of saying \emph{how} to traverse the data (loops, indexes, pointers), you simply say \emph{what} you want. The system finds the most efficient path. Fifty years later, SQL remains the universal language of relational databases, used by millions of developers every day.

\section{Introduction to SQL}

\begin{definition}[SQL]
\emph{SQL} (\emph{Structured Query Language}) is the standard language
for querying and manipulating relational databases. Unlike relational
algebra, SQL is a \emph{declarative} language: you specify \emph{what}
you want, not \emph{how} to get it.
\end{definition}

SQL has been standardized by ISO/ANSI. Major versions include:
SQL-86, SQL-92, SQL:1999, SQL:2003, SQL:2011, SQL:2016, SQL:2023.

\section{Structure of a SELECT Query}

\begin{keyformulas}[title=\textbf{General SELECT Syntax}]
\begin{minted}{sql}
SELECT [DISTINCT] columns
FROM   tables
[WHERE  condition]
[ORDER BY column [ASC|DESC]]
[LIMIT  n [OFFSET m]];
\end{minted}
\end{keyformulas}

The logical execution order differs from the writing order:
\begin{enumerate}
    \item \texttt{FROM} --- determine source tables.
    \item \texttt{WHERE} --- filter rows.
    \item \texttt{SELECT} --- project columns.
    \item \texttt{DISTINCT} --- eliminate duplicates.
    \item \texttt{ORDER BY} --- sort the result.
    \item \texttt{LIMIT / OFFSET} --- limit the number of rows.
\end{enumerate}

\section{SELECT and FROM}

\begin{codeblock}[title=\textbf{Selecting All Columns}]
\begin{minted}{sql}
SELECT * FROM Students;
\end{minted}
\end{codeblock}

\begin{codeoutput}
\begin{verbatim}
student_id | last_name | first_name | birth_date | major
-----------+-----------+------------+------------+------
1          | Brown     | Alice      | 2004-03-15 | CS
2          | Davis     | Bob        | 2003-11-22 | CS
3          | Miller    | Claire     | 2004-07-08 | Math
4          | Wilson    | David      | 2003-01-30 | CS
5          | Taylor    | Emma       | 2004-09-12 | Math
\end{verbatim}
\end{codeoutput}

\begin{codeblock}[title=\textbf{Selecting Specific Columns}]
\begin{minted}{sql}
SELECT last_name, first_name, major FROM Students;
\end{minted}
\end{codeblock}

\begin{bestpractice}[title=\textbf{Avoid SELECT *}]
In production, avoid \texttt{SELECT *}:
\begin{itemize}
    \item It transfers unnecessary data.
    \item Results depend on the schema: adding a column may break code.
    \item Use \texttt{SELECT *} only for exploration and debugging.
\end{itemize}
\end{bestpractice}

\section{Column and Table Aliases}

\begin{codeblock}[title=\textbf{Using Aliases}]
\begin{minted}{sql}
SELECT s.last_name AS student_name,
       s.first_name AS student_first,
       s.major AS department
FROM Students AS s;
\end{minted}
\end{codeblock}

The \texttt{AS} keyword is optional in most DBMS:
\texttt{SELECT last\_name student\_name FROM Students s}.

\section{WHERE Clause --- Filtering}

\begin{definition}[WHERE Clause]
The \texttt{WHERE} clause filters rows based on a Boolean condition.
Only rows for which the condition evaluates to \texttt{TRUE} are kept
(\texttt{FALSE} and \texttt{UNKNOWN} are rejected).
\end{definition}

\subsection{Comparison Operators}

\begin{keyformulas}[title=\textbf{SQL Comparison Operators}]
\begin{center}
\begin{tabular}{ll}
\toprule
\textbf{Operator} & \textbf{Meaning} \\
\midrule
\texttt{=} & Equal \\
\texttt{<>} or \texttt{!=} & Not equal \\
\texttt{<}, \texttt{>} & Less than, greater than \\
\texttt{<=}, \texttt{>=} & Less/greater than or equal \\
\texttt{BETWEEN a AND b} & Between $a$ and $b$ (inclusive) \\
\texttt{IN (v1, v2, ...)} & Member of a list \\
\texttt{LIKE 'pattern'} & Pattern matching \\
\texttt{IS NULL} / \texttt{IS NOT NULL} & Null test \\
\bottomrule
\end{tabular}
\end{center}
\end{keyformulas}

\begin{codeblock}[title=\textbf{Filtering Examples}]
\begin{minted}{sql}
-- CS students
SELECT * FROM Students WHERE major = 'CS';

-- Grades between 70 and 90
SELECT * FROM Enrollments WHERE grade BETWEEN 70 AND 90;

-- Names starting with 'B'
SELECT * FROM Students WHERE last_name LIKE 'B%';

-- Students in CS or Math
SELECT * FROM Students WHERE major IN ('CS', 'Math');

-- Enrollments with no grade
SELECT * FROM Enrollments WHERE grade IS NULL;
\end{minted}
\end{codeblock}

\subsection{Logical Operators}

\begin{codeblock}[title=\textbf{Combining Conditions}]
\begin{minted}{sql}
-- AND: both conditions must be true
SELECT * FROM Students
WHERE major = 'CS' AND birth_date > '2004-01-01';

-- OR: at least one condition true
SELECT * FROM Enrollments
WHERE grade > 90 OR grade < 50;

-- NOT: negation
SELECT * FROM Students
WHERE NOT major = 'Math';
\end{minted}
\end{codeblock}

\begin{warning}[title=\textbf{Three-Valued Logic and NULL}]
SQL uses \emph{three-valued logic}: \texttt{TRUE}, \texttt{FALSE}, \texttt{UNKNOWN}.
\begin{center}
\begin{tabular}{c|ccc}
\texttt{AND} & T & F & U \\
\hline
T & T & F & U \\
F & F & F & F \\
U & U & F & U \\
\end{tabular}
\qquad
\begin{tabular}{c|ccc}
\texttt{OR} & T & F & U \\
\hline
T & T & T & T \\
F & T & F & U \\
U & T & U & U \\
\end{tabular}
\end{center}
\texttt{NOT UNKNOWN = UNKNOWN}. Consequence: \texttt{grade = NULL}
returns \texttt{UNKNOWN}, not \texttt{TRUE}. Use \texttt{IS NULL}.
\end{warning}

\subsection{LIKE Patterns}

\begin{keyformulas}[title=\textbf{LIKE Wildcards}]
\begin{center}
\begin{tabular}{ll}
\toprule
\textbf{Wildcard} & \textbf{Meaning} \\
\midrule
\texttt{\%} & Zero or more arbitrary characters \\
\texttt{\_} & Exactly one arbitrary character \\
\bottomrule
\end{tabular}
\end{center}
\end{keyformulas}

\begin{codeblock}[title=\textbf{LIKE Examples}]
\begin{minted}{sql}
-- Names starting with 'M'
SELECT * FROM Students WHERE last_name LIKE 'M%';

-- First names with exactly 5 characters
SELECT * FROM Students WHERE first_name LIKE '_____';

-- Names containing 'il'
SELECT * FROM Students WHERE last_name LIKE '%il%';
\end{minted}
\end{codeblock}

\section{ORDER BY --- Sorting}

\begin{codeblock}[title=\textbf{Sorting Results}]
\begin{minted}{sql}
-- Ascending sort (default)
SELECT last_name, first_name FROM Students ORDER BY last_name ASC;

-- Descending sort
SELECT * FROM Enrollments ORDER BY grade DESC;

-- Multi-column sort
SELECT * FROM Students ORDER BY major ASC, last_name ASC;
\end{minted}
\end{codeblock}

\begin{codeoutput}
\begin{verbatim}
-- Sorted by grade descending:
student_id | course_id | grade | year
-----------+-----------+-------+-----
5          | 104       | 97.0  | 2025
3          | 102       | 95.5  | 2025
2          | 103       | 91.0  | 2025
1          | 101       | 88.5  | 2025
3          | 104       | 82.0  | 2025
1          | 102       | 76.0  | 2025
2          | 101       | 72.0  | 2025
4          | 101       | 55.0  | 2025
\end{verbatim}
\end{codeoutput}

\section{DISTINCT --- Eliminating Duplicates}

\begin{codeblock}[title=\textbf{Using DISTINCT}]
\begin{minted}{sql}
-- Without DISTINCT: duplicates possible
SELECT major FROM Students;
-- CS, CS, Math, CS, Math

-- With DISTINCT: unique values
SELECT DISTINCT major FROM Students;
-- CS, Math
\end{minted}
\end{codeblock}

\section{LIMIT and OFFSET --- Pagination}

\begin{codeblock}[title=\textbf{Limiting Results}]
\begin{minted}{sql}
-- Top 3 grades
SELECT s.last_name, e.grade
FROM Students s
JOIN Enrollments e ON s.student_id = e.student_id
ORDER BY e.grade DESC
LIMIT 3;

-- Pagination: page 2 (rows 4 to 6)
SELECT * FROM Students ORDER BY student_id LIMIT 3 OFFSET 3;
\end{minted}
\end{codeblock}

\begin{codeoutput}
\begin{verbatim}
-- Top 3 grades:
last_name | grade
----------+------
Taylor    | 97.0
Miller    | 95.5
Davis     | 91.0
\end{verbatim}
\end{codeoutput}

\begin{warning}[title=\textbf{LIMIT Is Not Standard SQL}]
\texttt{LIMIT} is supported by PostgreSQL, MySQL, SQLite.
The SQL:2008 standard uses \texttt{FETCH FIRST n ROWS ONLY}:
\begin{minted}{sql}
SELECT * FROM Students
ORDER BY last_name
FETCH FIRST 3 ROWS ONLY;  -- Standard SQL
\end{minted}
\end{warning}

\section{Expressions and Scalar Functions}

\begin{codeblock}[title=\textbf{Computed Expressions}]
\begin{minted}{sql}
-- Arithmetic expression
SELECT last_name, grade, ROUND(grade / 100.0 * 20, 1) AS grade_20
FROM Enrollments e
JOIN Students s ON e.student_id = s.student_id;

-- String concatenation
SELECT last_name || ' ' || first_name AS full_name
FROM Students;

-- Conditional expression CASE
SELECT last_name, grade,
    CASE
        WHEN grade >= 90 THEN 'Excellent'
        WHEN grade >= 80 THEN 'Very Good'
        WHEN grade >= 70 THEN 'Good'
        WHEN grade >= 60 THEN 'Satisfactory'
        ELSE 'Failing'
    END AS letter_grade
FROM Enrollments e
JOIN Students s ON e.student_id = s.student_id;
\end{minted}
\end{codeblock}

\begin{codeoutput}
\begin{verbatim}
last_name | grade | letter_grade
----------+-------+-------------
Brown     | 88.5  | Very Good
Brown     | 76.0  | Good
Davis     | 72.0  | Good
Davis     | 91.0  | Excellent
Miller    | 95.5  | Excellent
Miller    | 82.0  | Very Good
Wilson    | 55.0  | Failing
Taylor    | 97.0  | Excellent
\end{verbatim}
\end{codeoutput}

\section{Useful Functions}

\begin{keyformulas}[title=\textbf{Common Scalar Functions}]
\begin{center}
\begin{tabular}{lll}
\toprule
\textbf{Category} & \textbf{Function} & \textbf{Description} \\
\midrule
Text & \texttt{UPPER(s)}, \texttt{LOWER(s)} & Uppercase, lowercase \\
     & \texttt{LENGTH(s)} & Length \\
     & \texttt{SUBSTR(s, start, len)} & Substring \\
     & \texttt{TRIM(s)} & Remove whitespace \\
     & \texttt{REPLACE(s, a, b)} & Replace \\
Number & \texttt{ABS(x)}, \texttt{ROUND(x, n)} & Absolute value, rounding \\
       & \texttt{CEIL(x)}, \texttt{FLOOR(x)} & Ceiling, floor \\
Date   & \texttt{CURRENT\_DATE} & Current date \\
       & \texttt{DATE('now')} & Current date (SQLite) \\
NULL   & \texttt{COALESCE(a, b, c)} & First non-NULL value \\
       & \texttt{NULLIF(a, b)} & NULL if $a = b$ \\
\bottomrule
\end{tabular}
\end{center}
\end{keyformulas}

\begin{codeblock}[title=\textbf{Function Examples}]
\begin{minted}{sql}
SELECT UPPER(last_name) AS upper_name,
       LENGTH(first_name) AS name_len,
       COALESCE(grade, 0) AS grade_or_zero
FROM Students s
LEFT JOIN Enrollments e ON s.student_id = e.student_id;
\end{minted}
\end{codeblock}

\section{INSERT, UPDATE, DELETE}

\begin{codeblock}[title=\textbf{Inserting Data}]
\begin{minted}{sql}
-- Insert a single row
INSERT INTO Students (student_id, last_name, first_name, major)
VALUES (6, 'Garcia', 'Lucy', 'CS');

-- Insert multiple rows
INSERT INTO Students (student_id, last_name, first_name, major) VALUES
    (7, 'White', 'Mark', 'Math'),
    (8, 'Thomas', 'Julie', 'CS');

-- Insert from a query
INSERT INTO Archived_Students
SELECT * FROM Students WHERE major = 'Math';
\end{minted}
\end{codeblock}

\begin{codeblock}[title=\textbf{Updating and Deleting}]
\begin{minted}{sql}
-- Update
UPDATE Enrollments SET grade = 65.0
WHERE student_id = 4 AND course_id = 101;

-- Conditional delete
DELETE FROM Enrollments WHERE grade < 40;

-- Delete all rows
DELETE FROM Archived_Students;
-- or TRUNCATE TABLE Archived_Students;  (PostgreSQL)
\end{minted}
\end{codeblock}

\begin{warning}[title=\textbf{UPDATE/DELETE Without WHERE}]
An \texttt{UPDATE} or \texttt{DELETE} without a \texttt{WHERE} clause affects
\emph{all} rows in the table. Always verify your \texttt{WHERE} clause
before executing!
\end{warning}

\section{Python Integration}

\begin{codeblock}[title=\textbf{Complete CRUD with Python/sqlite3}]
\begin{minted}{python}
import sqlite3

conn = sqlite3.connect(':memory:')  # in-memory database
cur = conn.cursor()

# Create and insert
cur.executescript('''
    CREATE TABLE Students (
        id INTEGER PRIMARY KEY, last_name TEXT, first_name TEXT, major TEXT
    );
    INSERT INTO Students VALUES (1, 'Brown', 'Alice', 'CS');
    INSERT INTO Students VALUES (2, 'Davis', 'Bob', 'CS');
    INSERT INTO Students VALUES (3, 'Miller', 'Claire', 'Math');
''')

# Read with parameters
major = 'CS'
cur.execute("SELECT * FROM Students WHERE major = ?", (major,))
print("CS Students:", cur.fetchall())

# Update
cur.execute("UPDATE Students SET major = ? WHERE id = ?", ('Math', 2))
conn.commit()

# Verify
cur.execute("SELECT * FROM Students ORDER BY id")
for row in cur.fetchall():
    print(row)

conn.close()
\end{minted}
\end{codeblock}

\begin{codeoutput}
\begin{verbatim}
CS Students: [(1, 'Brown', 'Alice', 'CS'), (2, 'Davis', 'Bob', 'CS')]
(1, 'Brown', 'Alice', 'CS')
(2, 'Davis', 'Bob', 'Math')
(3, 'Miller', 'Claire', 'Math')
\end{verbatim}
\end{codeoutput}

\section*{Exercises}

\begin{exercise}
Write SQL queries for the following:
\begin{enumerate}
    \item All courses with more than 3 credits.
    \item Names of students born before 2004.
    \item Enrollments with a grade above 85, sorted by grade descending.
    \item The top 3 students by grade.
\end{enumerate}
\end{exercise}

\begin{exercise}
Explain the difference between the two queries:
\begin{minted}{sql}
SELECT DISTINCT major FROM Students;
SELECT major FROM Students GROUP BY major;
\end{minted}
\end{exercise}

\begin{exercise}
Write a SQL query that displays for each enrollment: the student's name,
the course title, the grade, and the corresponding letter grade
(Excellent $\geq$ 90, Very Good $\geq$ 80, Good $\geq$ 70,
Satisfactory $\geq$ 60, Failing otherwise).
\end{exercise}

\begin{exercise}
What are the results of the following expressions?
\begin{enumerate}
    \item \texttt{NULL = NULL}
    \item \texttt{NULL <> 5}
    \item \texttt{NULL AND TRUE}
    \item \texttt{NULL OR TRUE}
    \item \texttt{NOT NULL}
\end{enumerate}
\end{exercise}

\begin{exercise}
Write a Python program that:
\begin{enumerate}
    \item Creates a SQLite database with the University schema tables.
    \item Inserts the sample data.
    \item Executes 3 different queries and displays formatted results.
    \item Uses parameterized queries for all user inputs.
\end{enumerate}
\end{exercise}
